% Consolidated guide-ch13 from f8f4709b0ffac9da4c1fecbf2ce873002f6da6cc; source fragments preserved.

% === BEGIN guide-ch13-examples.tex blob c35545168f16e651dbdb66c94d5d4d935900bae0 ===
% 《大学物理学学习指导》第5版，第13章“光的衍射”，例13.1—13.8
% 书内印刷190—195页，用户源PDF物理198—203页。
% 解题原理、光栅缺级、角宽度等独立推演；次极大保留近似条件。

\sourceexample{13.1}{190}{198}
\problemtext 分析单缝衍射各级明条纹的角宽度，
并讨论缩窄/加宽狭缝、入射平面波相对主轴倾斜、
以及保持透镜不动仅上下移动狭缝时远场条纹的变化。
\workings
单缝宽$a$时暗纹满足$a\sin\theta_m=m\lambda$。
中央亮纹由$\pm1$级暗纹夹住，
\[
\boxed{\Delta\theta_{\rm center}
=2\arcsin(\lambda/a)\simeq2\lambda/a}.
\]
旁轴区相邻非中央亮纹宽度约为
$\boxed{\Delta\theta_{\rm side}\simeq\lambda/a}$。
$a$越小衍射条纹越宽、越明显；
$a\gg\lambda$时图样集中在几何光学焦点附近。
若准直入射方向相对主轴倾斜，中心亮纹随入射方向
移至透镜后焦平面的另一位置，整组衍射图样一并平移。
但在入射波面、焦面及薄透镜近似都不变的理想布置下，
仅\emph{平移孔径}不改变强度傅里叶变换的模，
所以不引起远场强度分布的平移。
\crosscheck 中央亮纹角宽约为相邻侧纹的两倍，
这一比值在小角近似下最清楚。
若$\lambda/a>1$，一阶暗纹无传播方向，中央亮纹覆盖可见角域，
上述普通“各级”划分不再适用。

\sourceexample{13.2}{191}{199}
\problemtext 单缝宽$a=0.60$ mm，透镜焦距$f=40$ cm，
焦面距中央亮纹$y=1.4$ mm处观察到某条侧明纹；
由可见波段$400$—$760$nm反推可能波长、
侧明纹级次与半波带数。
\workings
旁轴角度$\theta\simeq y/f$，单缝程差
\[
\delta=a\sin\theta\simeq a\frac yf
=\frac{(0.60\times10^{-3})(1.4\times10^{-3})}{0.40}
=2.1\times10^{-6}\,\mathrm m.
\]
用\emph{教材半波带近似}描述第$k$侧明纹：
$\delta\simeq(k+\tfrac12)\lambda$，有
\[
\lambda\simeq\frac{2100\,\mathrm{nm}}{k+1/2}.
\]
要求可见光，整数$k$只能为$k=3$或4；
分别得到
\[
\boxed{k=3:\ \lambda\simeq600\,\mathrm{nm},\ N_{\rm half}=7},
\]
\[
\boxed{k=4:\ \lambda\simeq466.7\,\mathrm{nm},\ N_{\rm half}=9}.
\]
\crosscheck \textbf{近似性质：}严格次极大位置需由
$\tan\beta=\beta$求非零根。
半波带奇数估算法给600nm、466.7nm，
并非严格极值处的精确波长；此题按原书近似方法整理。

\sourceexample{13.3}{191--192}{199--200}
\problemtext 双缝中心距$d=0.10$ mm，每条缝宽$a=0.020$ mm，
波长$\lambda=480$nm、透镜焦距$f=50$cm。
求单缝衍射中央亮纹宽度，以及中央包络内有多少条双缝干涉亮纹。
\workings
单缝第一级暗纹满足$\sin\theta_1\simeq\lambda/a$，
在焦面上半线宽$\ell=f\lambda/a$，
中央全宽
\[
\boxed{W=2f\lambda/a
=\frac{2(0.50)(480\times10^{-9})}{0.020\times10^{-3}}
=24\,\mathrm{mm}}.
\]
双缝干涉亮纹满足$d\sin\theta=m\lambda$。
在中央包络内部$|\sin\theta|<\lambda/a$，
故$|m|<d/a=5$。
$m=\pm5$处恰遇单缝第一暗纹，属于缺级，因此
\[
\boxed{m=0,\pm1,\pm2,\pm3,\pm4},\qquad
\boxed{N_{\rm visible}=9}.
\]
\crosscheck 缺级导致中央包络内亮纹是9条，不是11条。

\sourceexample{13.4}{192--193}{200--201}
\problemtext 光栅垂直入射$\lambda=600$nm时，
第2、3级亮纹分别位于$\sin\theta_2=0.20,\sin\theta_3=0.30$，
第4级光栅主极大缺级。
求光栅常量$d=a+b$、可能的狭缝宽度$a$、
符合题意的最小$a$及最终全部可见主明纹级数。
\workings
光栅方程$d\sin\theta_k=k\lambda$，所以
\[
\boxed{d=\frac{2(600\,\mathrm{nm})}{0.20}
=6.0\,\mu\mathrm m}.
\]
缺第4级意味着光栅主极大与单缝某阶零点重合：
$d/a=4/m$，其中$m=1,2,3$且$a<d$。
因此有
\[
\boxed{a=1.5,\ 3.0,\ 4.5\,\mu\mathrm m}.
\]
若按原题要求取最小值$a=1.5\,\mu$m，
$d/a=4$，第4、8等4的倍数级缺失。
又可见传播方向须满足$|k|<d/\lambda=10$，
严格排除掠射$90^\circ$的第10级，
\[
\boxed{k=0,\pm1,\pm2,\pm3,\pm5,\pm6,\pm7,\pm9}
\]
共$\boxed{15}$条主亮纹。
\crosscheck 对另外两种允许缝宽，缺级周期不同；
“15条”特指采用源解所取\emph{最小}的狭缝宽度方案。

\sourceexample{13.5}{193}{201}
\problemtext 三块光栅的刻线密度分别为1400、600、100条/mm，
拟测量可见光400—700nm范围，选择适合观察分立衍射光谱的光栅。
\workings
对应光栅常量约为
\[
d_1=0.714\,\mu\mathrm m,\qquad
d_2=1.667\,\mu\mathrm m,\qquad
d_3=10.0\,\mu\mathrm m.
\]
第一级光栅方程$\sin\theta=\lambda/d$。
对于第1块在$\lambda=700$nm时，
$\sin\theta\simeq0.98$，第一级极接近掠射，
不便于实际布置与成像。
对第3块，一级角度约$2.3^\circ$—$4.0^\circ$，
谱线角分离较小。
第2块使一级波段约在
$\arcsin(0.24)$至$\arcsin(0.42)$，
约$14^\circ$—$25^\circ$，适于观测。
故按题设选\textbf{600条/mm的第二块光栅}。
\crosscheck 刻线密度越高不保证量程越好；
还应同时考虑有限传播角、光谱重叠与仪器几何范围。

\sourceexample{13.6}{193--194}{201--202}
\problemtext 一平面光栅每厘米刻590条，垂直入射波段为
$0.4$—$0.7\,\mu$m的平行光。
求第一级谱的角宽度；判断第3级紫光与第2级红光
是否有重叠；若要求第二级全部落在单缝中央衍射包络内，
估算狭缝宽度的最大允许值。
\workings
光栅常量
\[
d=\frac{1\,\mathrm{cm}}{590}
\simeq1.695\times10^{-5}\,\mathrm m.
\]
第一级两端角
$\theta_V=\arcsin(400\,\mathrm{nm}/d)\simeq1.35^\circ$，
$\theta_R=\arcsin(700\,\mathrm{nm}/d)\simeq2.37^\circ$，
因而
\[
\boxed{\Delta\theta_1\simeq1.02^\circ}.
\]
第3级400nm光的光程级数积为1200nm，
第2级700nm光的级数积为1400nm，
对应不同角度，但两级\emph{连续波段}在
$3\lambda_V\le2\lambda_R$时有交叠，
故\textbf{第3级与第2级光谱存在波段重叠}。
要让第二级红端仍处于单缝中央包络内，
第一暗纹界限$a\sin\theta=\lambda$应不小于
二级红光衍射角，即
\[
a\lesssim\frac{\lambda_R}{\sin\theta_{2,R}}
=\frac{\lambda_R}{2\lambda_R/d}
=\frac d2
\simeq8.47\times10^{-6}\,\mathrm m.
\]
故$\boxed{a_{\max}\simeq8.5\,\mu\mathrm m}$。
\crosscheck \textbf{原书题面与印刷参考解存在单位矛盾：}
题面确实印为\emph{每厘米590条}，即上述$d\simeq16.95\,\mu$m，
角宽约$1.015^\circ$、$a_{\max}\simeq8.47\,\mu$m。
然而原书参考解改用$d\simeq1.69\times10^{-4}$ cm
（即$1.69\,\mu$m、对应\emph{每毫米590条}），
并给出首级光谱角宽约$10^\circ48'$、
$a_{\max}\simeq0.85\,\mu$m。
两套数值不能同时满足题面刻痕密度；
本稿尊重题面字面解，另外保留原书答案所采用的单位假设，
此题\emph{原书单位勘误待二校}。


\sourceexample{13.7}{194--195}{202--203}
\problemtext 两波长$\lambda_1=440$nm、$\lambda_2=660$nm垂直入射同一光栅，
两者非零主明纹的\emph{第二次重合}发生在衍射角$60^\circ$，
求光栅常量。
\workings
两束光在同一衍射方向重合须满足
$k_1\lambda_1=k_2\lambda_2$，
故
\[
\frac{k_1}{k_2}=\frac{660}{440}=\frac32.
\]
最小正整数重合是$(k_1,k_2)=(3,2)$，
第二次非零重合为$(6,4)$。
于是光栅方程
$d\sin60^\circ=6(440\,\mathrm{nm})$，
\[
\boxed{d=\frac{2640\,\mathrm{nm}}{\sin60^\circ}
\simeq3048\,\mathrm{nm}\simeq3.05\,\mu\mathrm m}.
\]
\crosscheck 这里“第二次重合”不包括两色均重合的零级中央主明纹；
若把中央零级也计入，则整数级次会不同。

\sourceexample{13.8}{195}{203}
\problemtext 晶体晶面距$d=0.275$nm，含有波长
$0.095$—$0.130$nm的X射线束以$45^\circ$的晶面夹角入射，
求哪些波长能产生强烈布拉格反射。
\workings
由布拉格定律
$2d\sin\theta=k\lambda$，$\theta=45^\circ$，
\[
\lambda_k=\frac{2(0.275)\sin45^\circ}{k}\,\mathrm{nm}
=\frac{0.38891}{k}\,\mathrm{nm}.
\]
逐个检查正整数$k$：
\[
\begin{aligned}
k=1&:\ 0.389\,\mathrm{nm},\\
k=2&:\ 0.194\,\mathrm{nm},\\
k=3&:\ 0.1296\,\mathrm{nm},\\
k=4&:\ 0.0972\,\mathrm{nm},\\
k=5&:\ 0.0778\,\mathrm{nm}.
\end{aligned}
\]
落在原题给定范围的只有
\[
\boxed{\lambda_3\simeq0.130\,\mathrm{nm}},\qquad
\boxed{\lambda_4\simeq0.097\,\mathrm{nm}}.
\]
\crosscheck $45^\circ$是射线与晶面的夹角（布拉格掠射角），
不是与晶面法线的夹角；在本题恰为45度所以两种角度数值相同，
但一般题不能混用。
% === END guide-ch13-examples.tex ===
